\chapter{Problem Definition and Requirements}
\label{chap:requirements}

\section{Problem Statement}
\label{sec:req-problem}

The volume of AI-related material published daily (research preprints,
model and framework releases, government and policy notices, funding
announcements, long-form video, and discussion on aggregator platforms) has
grown past what any individual reader can triage manually across the many
independent sites, feeds, and channels that publish it. Addressing this
information-overload problem requires more than mechanically republishing
content: the system must continuously acquire material from a heterogeneous
and growing set of source types, reduce each item once to structured
information later processes can reason about without re-reading raw text,
remove redundancy created by multiple outlets covering the same event, rank
and filter what remains against an individual user's demonstrated interests,
and expose both a browsable feed and a natural-language interface for
questions a static ranked list cannot answer. Each of these needs is
restated below as a concrete requirement, realized by a component described
in the chapters that follow.

\section{Functional Requirements}
\label{sec:req-functional}

\begin{description}
  \item[FR1: Multi-source ingestion via a pluggable adapter pattern.]
  New source types must be addable without modifying the agents and
  services that consume their output, since the set of sources this system
  covers is expected to keep growing.

  \item[FR2: Automated relevance screening for user-submitted sources.]
  Since end users, not only maintainers, may register new sources, the
  system must accept or reject a candidate source's topical fit without a
  human reviewer, so the source catalog can scale with the user base.

  \item[FR3: Single-pass content enrichment.]
  Every item must be reduced once to a structured record (summary, topics,
  entities, category) so ranking, search, clustering, and the assistant can
  query stored facts instead of re-deriving them from raw text on every
  read.

  \item[FR4: Near-duplicate clustering.]
  Because independent outlets frequently cover the same story, the system
  must detect and group near-duplicate items so a user is not shown the
  same event repeatedly under different bylines.

  \item[FR5: Personalized, explainable, deterministic ranking.]
  The system must order each user's feed by that user's demonstrated
  interests and attach a human-readable reason to every ranked item, so the
  ordering is auditable rather than an opaque score.

  \item[FR6: Semantic search.]
  Users must be able to locate content by meaning rather than exact
  keyword match, requiring retrieval over a vector representation of each
  item.

  \item[FR7: A grounded conversational assistant.]
  Users must be able to ask free-form questions about the corpus or a
  specific item and receive an answer traceable to specific source
  passages, since an ungrounded answer about news content cannot be
  trusted without a way to verify it.

  \item[FR8: Deep-media transcription and chaptering.]
  Long-form video lacking a usable caption track must still enter the same
  enrichment and ranking pipeline as text, requiring a transcription
  fallback and, for sufficiently long videos, chapter-level summarization
  rather than one pass over the whole transcript.

  \item[FR9: Free/Pro entitlements and billing.]
  The system must gate a subset of features behind a paid plan and process
  real payments for it, since a system meant to operate as a product needs
  a mechanism to sustain its own running costs.
\end{description}

\section{Non-Functional Requirements}
\label{sec:req-nonfunctional}

\begin{description}
  \item[Determinism in the ranking serving path.]
  The function that orders a feed at request time must not invoke a large
  language model, since an LLM call is comparatively slow, costly, and
  non-reproducible for an operation performed on every request.

  \item[Per-ORM database ownership boundaries.]
  With two ORMs reading and writing one database, each table must be owned
  by exactly one ORM, which alone creates and migrates it; the other ORM
  may only read it through a dedicated mirror model, never write or
  migrate across the boundary.

  \item[Idempotent ingestion.]
  Re-running an ingestion phase, after a crash, retry, or scheduling
  overlap, must not create duplicate rows, since the pipeline runs
  unattended and cannot rely on manual cleanup.

  \item[Graceful degradation when an optional integration is unconfigured.]
  Billing, transactional email, and LLM access depend on third-party
  services that may be absent in a given deployment; the system must keep
  functioning with those features inactive rather than failing outright.

  \item[Citation-verifiable RAG answers.]
  Every claim the assistant attributes to a source must be checkable
  against that source at generation time; an answer whose citations cannot
  be resolved must be flagged as such, not presented as equally
  trustworthy.

  \item[Reproducible evaluation.]
  Claims about ranking quality must come from a versioned, re-runnable
  evaluation procedure over stored data, not one-off manual inspection, so
  a later run on the same data yields the same verdict.
\end{description}

\section{Architecture Principles}
\label{sec:req-principles}

The requirements above are enforced, in this project, by a small set of
architecture principles adopted early and treated as invariants throughout
development. Four recur across the chapters that follow and are stated here
as formal design constraints:

\begin{enumerate}
  \item \textbf{Ingest once, personalize later.} Content is scraped,
  enriched, embedded, scored, and clustered exactly once, globally; per-user
  work happens only at read/serve time over that shared pool. Scraping or
  enriching content separately per user is disallowed.
  \item \textbf{One enrichment call per item.} All structured metadata for
  an item (summary, topics, entities, category) comes from a single
  language-model call at ingest time; a new metadata field extends that
  call's output schema, never adds a second pass over the corpus.
  \item \textbf{Recommendation trends away from the language model in the
  hot path.} A language model's role in ranking is limited to producing a
  human-readable explanation for a score already computed deterministically;
  the score itself is never computed by a language model at serving time.
  \item \textbf{Strict per-ORM table ownership.} Every table has exactly one
  owning ORM; cross-ORM communication is read-only in both directions, and a
  feature that appears to need a cross-ORM write is instead resolved by
  adding a table on the writer's own side and mirroring it, never by a
  one-off write exception.
\end{enumerate}

Chapter~\ref{chap:architecture} shows how these principles are realized in
the system's codebase structure, and later chapters revisit each principle
where it governs a specific subsystem's design.
